\recipeTitle{Ciambellone all'arancia e latte di mandorla}

\recipeSubtitle{Un ciambellone profumato all'arancia, soffice e leggermente aromatico grazie al latte di mandorle tostate.}

\recipeAuthor{Felice Pantaleo}

\recipeStory{Una ricetta casalinga nata attorno a una planetaria, un'arancia biologica e il sacro diritto di assaggiare l'impasto rimasto nella ciotola.}

% Add a real image in recipes/ciambellone-arancia-latte-mandorla/images/ and uncomment the line below.
% \recipePhoto{recipes/ciambellone-arancia-latte-mandorla/images/ciambellone-arancia-latte-mandorla.jpg}

\recipeInfo
  {20 minuti}
  {40 minuti}
  {1 ciambellone}
  {Facile}
  {Dolce da colazione}

\begin{ingredients}
  \ingredient{Farina per dolci}{250 g}
  \ingredient{Zucchero a velo}{250 g, piu' q.b. per decorare}
  \ingredient{Uova grosse fresche}{4}
  \ingredient{Lievito Pane Angeli}{1 bustina}
  \ingredient{Arancia biologica}{1}
  \ingredient{Burro vegetale}{180 g}
  \ingredient{Latte di mandorle tostate senza zuccheri aggiunti}{50 g}
  \ingredient{Burro vegetale e farina}{q.b. per lo stampo}
\end{ingredients}

\begin{tools}
  \tool{Planetaria con frusta}
  \tool{Ciotola}
  \tool{Setaccio}
  \tool{Grattugia fine}
  \tool{Spremiagrumi}
  \tool{Teglia per ciambellone}
  \tool{Spatola}
\end{tools}

\begin{procedure}
  \step{Preriscalda il forno a 170\textdegree C in modalita' ventilata, con calore sia dal basso sia dall'alto.}

  \step{Lava bene 1 arancia biologica, grattugia finemente la scorza evitando la parte bianca, poi spremi il succo ed elimina eventuali noccioli.}

  \step{Nella planetaria unisci 4 uova grosse fresche e 250 g di zucchero a velo. Monta con la frusta a velocita' sostenuta per circa 2 minuti, finche' il composto diventa chiaro, uniforme e leggermente spumoso.}

  \step{Aggiungi 180 g di burro vegetale poco alla volta, un cucchiaino per volta, continuando a montare. Aspetta che ogni aggiunta sia parzialmente incorporata prima di aggiungere la successiva.}

  \step{Dopo circa 2 minuti, aggiungi la scorza grattugiata di 1 arancia biologica e il suo succo spremuto. Continua a montare fino a incorporarli bene nell'impasto.}

  \step{Setaccia 250 g di farina per dolci. Aggiungila poco alla volta nella planetaria, continuando a lavorare l'impasto a velocita' media, quanto basta per ottenere una massa liscia e senza grumi.}

  \step{Versa a filo 50 g di latte di mandorle tostate senza zuccheri aggiunti, continuando a mescolare.}

  \step{Aggiungi 1 bustina di lievito Pane Angeli e monta ancora per circa 2 minuti, in modo da distribuire bene il lievito nell'impasto.}

  \step{Imburra con burro vegetale e infarina accuratamente le pareti interne di una teglia per ciambellone. Versa l'impasto nello stampo e livella leggermente la superficie con una spatola.}

  \step{Raccogli con il dito l'impasto rimasto nella ciotola. Lecca il dito: e' parte ufficiale del procedimento.}

  \step{Cuoci il ciambellone a 170\textdegree C in forno ventilato per circa 40 minuti. Controlla la cottura con uno stecchino: deve uscire asciutto o con poche briciole umide, ma senza impasto crudo.}

  \step{Sforna il ciambellone e lascialo raffreddare nello stampo per circa 15 minuti. Poi coprilo leggermente per non perdere troppa umidita'.}

  \step{Quando il ciambellone e' tiepido, staccalo delicatamente dalla teglia, trasferiscilo su un piatto da portata e setaccia sopra poco zucchero a velo.}
\end{procedure}

\begin{notes}
  \note{Il burro vegetale puo' essere aggiornato con marca e dettagli quando disponibili.}
  \note{La copertura durante il raffreddamento aiuta a mantenere il dolce piu' umido e soffice.}
\end{notes}
